% Included chapter fragment. Requires only the packages in the main preamble.
% This file deliberately defines no document class, fonts, or global macros.

\section{阅读地图：把文章当作一条可验证的优化链}

George Kopanas 的文章 \emph{What the 3DGS Literature Left on the Table}
不是一份 CUDA 技巧清单，而是一堂关于\textbf{瓶颈迁移}的课程：先消除不必要的
相交对，再降低原子操作和归约成本，然后处理寄存器、无用反向工作、分歧以及
执行流水线。前一个步骤改变了后一个步骤成立的条件。例如，没有精确的子块筛选，
把逐像素分支全部改成无分支计算，可能只是让更多无效像素消耗算力。

本章依次讲解 \texttt{s00} 至 \texttt{s06}。每一步都区分三件事：文章的设计、
本仓库所用 gsplat 分支原来具备的能力，以及这次移植实验实际保留的内容。
代码片段是解释依赖关系的教学伪代码，不是可以直接替换通用内核的补丁；
省略的边界、同步和张量布局检查不能在生产代码中省略。

\paragraph{先划清数字的适用边界。}
文章使用一张 NVIDIA A100-SXM4-40GB，在九个 Mip-NeRF 360 场景的
7k、30k 检查点上测量。其标题中的 $7.30\times$（30k）和 $8.38\times$（7k）
对应最终 BF16 SH 版本的前向加反向渲染流水线，与 Inria 参考实现比较；
它们不是任意训练器的优化器、损失、数据加载与增密循环的整体倍率。
文章最终 FP32 SH 版本在相应集合上的数字也不同：分别为 $6.54\times$ 和
$7.55\times$。不能把 BF16 的标题归到纯 FP32 版本上。

本任务使用 RTX 4090、20,789,814 个高斯、一阶 SH、$1280\times720$ 图像，
对比的是已经具备部分优化的 gsplat 分支。硬件、场景、初始实现、SH 阶数和
接口契约都不相同。因此\textbf{文章的 A100 计数器、启动参数和标题倍率均不可直接迁移}。
本章中有关本任务的结论以
\texttt{benchmarks/blog\_eval/REPORT.md} 的最终重建记录为准，
不把较早的探索性数字拼接成一个不存在的最终结果。

\subsection{先建立一个关于成本的模型}

记总高斯数为 $N$，当前相机可见的高斯数为 $M$，发出的 tile--Gaussian
相交对数为 $P$，实际执行的像素--Gaussian 混合次数为 $E$。
这四个数决定了不同阶段的主要工作量：
\begin{equation}
 t_{\mathrm{step}}\approx
 t_{\mathrm{project}}(N)+t_{\mathrm{SH}}(M)
 +t_{\mathrm{bin/sort}}(P,M)+t_{\mathrm{raster}}(E)
 +t_{\mathrm{VJP}}(N,M).
\end{equation}
这里 VJP 指向量--雅可比乘积，即反向传播中的链式法则计算。该式不是拟合出的
性能公式，而是一张诊断地图：如果排序占主导，优化归约并不会让排序自动变快；
如果反向预处理仍扫描全部 $N$，减少像素混合次数也不会消除所有大场景开销。

假设某阶段占原总时间的比例为 $f$，该阶段获得倍率 $s$，则整个调用最多变为
\begin{equation}
 S_{\mathrm{total}}=\frac{1}{(1-f)+f/s}.
\end{equation}
这解释了为什么本任务可以看到明显的 raster 加速，却只有较小的完整渲染收益。
还要注意，CUDA-event 的内核阶段之和与同步的端到端延迟不是同一个统计量：
启动、框架调度、等待和内存初始化可以出现在后者中。文章的一视图阶段图、
九场景均值，以及本任务三视图测量，不能混在同一张倍率表里。

\section{\texttt{s00}：忠实基线比一个慢基线更重要}

\subsection{五阶段流水线和不可随意改变的语义}

文章的 \texttt{s00\_baseline} 是在统一测试框架中忠实重写 Inria CUDA
光栅器，而不是故意制造一个低效版本。后续包逐步增加优化，使每一段差异可以
单独归因。文章对 Inria 的倍率仍以实测的 Inria 为分母，并不偷偷换成更慢的重写版本。

一个高斯的世界空间参数为位置、尺度、旋转、不透明度和 SH 系数。投影之后，
光栅器主要使用二维中心 $\boldsymbol\mu'$、逆协方差
\begin{equation}
 \Lambda=\begin{bmatrix}a&b\\b&c\end{bmatrix},\qquad
 q(\boldsymbol d)=a d_x^2+2b d_xd_y+c d_y^2,
 \qquad \boldsymbol d=\boldsymbol p-\boldsymbol\mu'.
\end{equation}
对正定的 $\Lambda$，$q\geq0$。高斯权重和混合不透明度为
\begin{equation}
 G=\exp(-q/2),\qquad
 \alpha_{\mathrm{raw}}=oG,\qquad
 \alpha=\min(\alpha_{\max},\alpha_{\mathrm{raw}}).
\end{equation}
这里只把 $o$ 定义为已经激活后的不透明度；若保存的是 logit，调用前需要 sigmoid。
同理，保存的对数尺度必须先取指数。把训练文件中的原始参数直接传给要求物理尺度的
接口，虽然可能仍能渲染，却已经不是同一个场景。本任务早期一次探索就犯过对数尺度
未激活的错误，那些时间和图像比较均已作废。

按深度从前到后，颜色与透射率的递推为
\begin{align}
 \boldsymbol C_k&=\boldsymbol C_{k-1}
                  +T_{k-1}\alpha_k\boldsymbol c_k,\\
 T_k&=T_{k-1}(1-\alpha_k),\qquad T_0=1,\\
 \boldsymbol I&=\boldsymbol C_K+T_K\boldsymbol c_{\mathrm{bg}}.
\end{align}
低于 cutoff 的 splat 不参与混合；像素达到透射率终止条件后停止。特别要区分：
检测终止的是更新前还是更新后的 $T$，触发终止的 splat 是先混合还是先拒绝，
以及停止索引采用包含还是不包含端点。这些是基线语义，不是实现小节。

\begin{lstlisting}[language=C++,basicstyle=\ttfamily\small,breaklines=true]
// Semantic skeleton: keep the reference's precise stop ordering.
project_and_shade();
count_pairs_then_prefix_scan();
emit_and_depth_sort_per_tile();
for (splat : front_to_back_list) {
    alpha = evaluate_with_reference_cutoff_and_cap(splat, pixel);
    if (!reference_accepts(alpha, T)) continue;
    C += (T * alpha) * color;
    T *= (1.0f - alpha);
    record_reference_last_index();
}
// Backward replays the accepted prefix in reverse order.
for (splat : reverse_accepted_prefix) {
    T_before = T_after / (1.0f - alpha);
    accumulate_screen_space_vjp();
}
apply_projection_covariance_and_SH_vjp();
\end{lstlisting}

\subsection{反向为什么昂贵}

一个高斯会影响多个像素、多个 warp、甚至多个 tile。屏幕空间梯度通常包含中心
2 项、conic 3 项、不透明度 1 项、RGB 3 项，共 9 项。若每个像素直接向全局数组
执行 9 次原子加，许多线程会竞争同一高斯地址。反向还需要逆序重建透射率：
\begin{equation}
 T_{k-1}=\frac{T_k}{1-\alpha_k}.
\end{equation}
因此 alpha cap 不只控制图像，也保证分母不为零；改变 cap 会同时改变渲染和梯度。
文章示意路径使用 $0.99$，本任务必须保留 gsplat 的 $0.999$。两者不能因为
``都很接近 1''就互换。

\paragraph{预期瓶颈与风险。}
忠实基线常见的瓶颈是宽松分箱产生过多 $P$、逐像素原子操作以及反向重放。
但是基线还承担可验证性职责：相等深度顺序、像素中心、背景梯度、clamp 导数、
空输入和异步流都必须固定。PSNR 相同不意味着参数梯度相同，图像误差很小也不能
单独证明反向正确。

\paragraph{本分支与本任务。}
本任务没有把 gsplat 改写成文章的 \texttt{s00}，而是以固定提交
\texttt{ffe2db3} 为自身基线。最终比较还固定参数激活、四元数归一化和 compact-box
配置。文章对 dense reference 的阈值提供了方法论，但不能代替本分支的差分测试。

\section{\texttt{s01}：先把五类已有文献优化讲透}

文章把 DISTWAR、Taming 3DGS、StopThePop、FlashGS、Speedy-Splat、
Splatshop、Faster-GS、HiGS、gsplat、LichtFeld Studio 等工作中的技术综合起来。
这一步的价值在于消除基线遗漏，不应将其全部算作本任务的新发明。

\subsection{第一类：warp butterfly 归约，先少写再谈快算}

同一个 warp 的 32 个像素线程在同一次迭代处理同一个高斯，设线程 $\ell$ 对
梯度分量 $j$ 的贡献为 $g_{\ell j}$。需要写入的是
\begin{equation}
 G_j=\sum_{\ell=0}^{31}g_{\ell j},\qquad j=0,\ldots,8.
\end{equation}
先在寄存器之间归约，最后仅由 lane 0 写每一项，可把同一 warp 的原子更新数
最多降低 32 倍。shuffle 是 lane 间寄存器交换，不经过共享内存。

\begin{lstlisting}[language=C++,basicstyle=\ttfamily\small,breaklines=true]
// Every lane must participate; rejected pixels supply exact zero.
for (int j = 0; j < 9; ++j) {
    float sum = grad[j];
    for (int d = 16; d != 0; d >>= 1)
        sum += __shfl_xor_sync(0xffffffffu, sum, d);
    if (lane == 0) atomicAdd(dst + j, sum);
}
\end{lstlisting}

5 次交换之后每个线程都有总和，但其中 31 份会被丢弃。9 个分量需要
$9\times5=45$ 条 shuffle，lane 0 还要发出 9 条原子指令。这里的
``一条指令''是 warp 级发射，不是一个地址上的交易量；减少指令数不等于消除跨 tile
对同一高斯的竞争。

\paragraph{瓶颈、风险和状态。}
该方法针对原子压力，随后瓶颈会迁移到 shuffle 和指令发射。必须保证归约成员都
处理同一个高斯，且不能让部分 lane 提前退出而仍使用完整掩码。本 fork 原来已有
逐分量全 warp 归约与 lane-0 原子写回；本任务的新内容是下一节的 reduce-scatter，
不是再次发明 warp 聚合。

\subsection{第二类：不透明度感知的精确行跨度}

原始 $3\sigma$ 外接圆再包成正方形，忽略了椭圆各向异性，也忽略了低不透明度
高斯的有效面积。令 cutoff 为 $\tau=1/255$，只要 $o\geq\tau$，有效区域满足
\begin{equation}
 o\exp(-q/2)\geq\tau
 \quad\Longleftrightarrow\quad
 q\leq\rho^2,\qquad \rho^2=2\ln(o/\tau).
\end{equation}
若 $o<\tau$，整个高斯都不能通过该 cutoff；若恰好相等，边界退化情形仍需按
比较符号处理。令 $D=ac-b^2>0$，椭圆的紧 AABB 半宽为
\begin{equation}
 x_{\max}=\rho\sqrt{c/D},\qquad
 y_{\max}=\rho\sqrt{a/D}.
\end{equation}
AABB 已经优于最大特征值圆，但倾斜椭圆的四角仍会产生无效 tile 测试。

对固定相对纵坐标 $y$，解二次方程可以得到横向边界
\begin{equation}
 x_\pm(y)=-\frac{b}{a}y
 \ \pm\ \sqrt{\frac{\rho^2}{a}-\frac{D}{a^2}y^2}.
\end{equation}
真正的 tile 行不是一条水平线，而是一个纵向带状区间 $[y_0,y_1]$。
该行跨度要取区间中 $x_-$ 的最小值与 $x_+$ 的最大值：检查区间端点以及落在
区间内的驻点，再将连续坐标映射成离散 tile 索引。只在行中心代入公式会漏掉
擦过上下边缘的椭圆；那不是精确行跨度。

\begin{lstlisting}[language=C++,basicstyle=\ttfamily\small,breaklines=true]
if (opacity < cutoff) return no_pairs;
rho2 = 2 * log(opacity / cutoff);
for (int ty = first_row; ty < end_row; ++ty) {
    Strip strip = reference_tile_strip(ty);
    Bounds x = extrema_of_ellipse_over_strip(conic, rho2, strip);
    Span span = conservative_discrete_tile_span(x, image_bounds);
    for (int tx = span.begin; tx < span.end; ++tx)
        emit(tx, ty, gaussian_id);
}
\end{lstlisting}

\paragraph{瓶颈与风险。}
这是减少 $P$ 的算法性优化，同时减轻发射、排序、共享 staging 与像素测试。
它往往比内层少一条乘法更值得先做。风险在于病态 conic、根号的小负舍入、
切线相交、像素中心约定和整数端点。浮点实现应保守，不允许为了少发一个 tile
而丢掉任何参考实现会接受的像素；对非正定或溢出情况应有明确回退。

\paragraph{本分支与本任务。}
本 fork 的 \texttt{isect\_tiles.cu} 已有 opacity-aware compact-box 和解析 sweep，
\texttt{compact\_box\_impl="sweep"} 是实现选择的默认值，但是否启用 compact-box
仍是另一个开关。产品封装原来已经默认启用。相机 0 的消融显示相交对约从
16.05M 降为 9.52M。这是已有能力的验证，不是本任务新增的收益。

\subsection{第三类：$8\times4$ 子块掩码和 48 字节结构体}

一个 $16\times16$ tile 有 256 个像素。让 8 个 warp 分别拥有一个
$8\times4$ 矩形，可把每个高斯对 8 个子块的几何重叠关系压成一个 8 位掩码。
对于 warp $w$，掩码的第 $w$ 位为 0，意味着整个 warp 可以完全跳过该高斯。
这比 32 个线程先分别计算高斯再投票更早地消除了工作。

把 32 个 staging 高斯看作一个小批次：lane $\ell$ 检查第 $\ell$ 个高斯是否
覆盖当前 warp，ballot 把 32 个布尔值组成工作表，再用最低置位迭代。

\begin{lstlisting}[language=C++,basicstyle=\ttfamily\small,breaklines=true]
unsigned work = __ballot_sync(0xffffffffu,
    lane < batch_count && ((batch_mask[lane] >> warp_id) & 1));
while (work != 0) {
    int j = __ffs(work) - 1;
    work &= work - 1;  // preserve increasing splat/depth order
    blend_shared_splat(j);
}

struct alignas(16) Splat48 {
    float4 center_and_aux;
    float4 conic_and_opacity;
    float4 rgb_and_aux;
};  // Illustrative field assignment: 3 * 16 bytes.
\end{lstlisting}

48 字节结构体将二维中心、conic、颜色和不透明度连续存储，使单线程可以用
3 次 16 字节宽读取完成 staging。此处的收益是减少散乱的标量地址与加载指令，
不是宣称每个 warp 只产生 48 字节流量。相邻线程访问不同高斯时，实际事务仍取决于
布局、缓存扇区和访问模式；\texttt{float4} 名称本身也不能保证最终机器码一定向量化。

\paragraph{瓶颈与风险。}
掩码针对 warp 级空转，打包针对 staging 指令和局部性。计算 mask 本身也有成本，
必须由跳过的工作偿还。布局改变还涉及 packed camera、多通道特征、对齐和 ABI。
几何 mask 只是``可能覆盖''，不是``确实贡献''，不能拿来直接证明该高斯的梯度非零。

\paragraph{本分支与本任务。}
本任务没有移植文章的精确 conic 子块 mask，也没有交付 48 字节打包结构。
曾测试 $8\times4$ 像素所有权加保守 AABB mask 的前向原型，FP32 版本约有
$1.020\times$ 的 raster 收益，但小且有噪声，并增加专用循环和回退，最终撤回。
这不是对文章精确算法的反证：AABB 更松，病态 conic 回退更多，原型并不等价。

\subsection{第四类：两阶段排序，把重复的深度工作提到前面}

传统方法为每个相交对构造 $({\rm tile},{\rm depth})$ 宽键，在 $P$ 个元素上排序。
同一个高斯的深度会随其覆盖的 tile 数重复参与排序。若先把 $M$ 个可见高斯按
32 位深度排序，再按此顺序发射 tile 对，整个相交数组已经具有全局深度顺序。
最后对 tile ID 做\textbf{稳定}排序，就能在按 tile 分组的同时保留组内深度顺序。

\begin{equation}
 \operatorname{stableSort}_{\rm tile}
 \left(\operatorname{emit}\left(
 \operatorname{sort}_{\rm depth}(\text{visible Gaussians})\right)\right)
 \ \Longrightarrow\ \text{每个 tile 内深度有序}.
\end{equation}

\begin{lstlisting}[language=C++,basicstyle=\ttfamily\small,breaklines=true]
visible = filter(count_tiles > 0);
visible = stable_sort_by_depth(visible);
offsets = exclusive_scan(count_tiles[visible]);
emit_at_scanned_offsets(visible, offsets, tile_keys, gaussian_ids);
stable_radix_sort(tile_keys, gaussian_ids);  // narrow tile key
find_tile_ranges();
\end{lstlisting}

扫描偏移很重要：如果每个高斯用无序的全局原子 append 发射，``深度顺序发射''
就不再成立。相等深度也需要沿用基线 tie-break；透明混合并不交换，稳定性是
图像语义的一部分。正的有限相机深度通常可以用相应浮点位序构造单调键；
负深度、NaN 或不同编码不能靠强制类型转换草率处理。

\paragraph{瓶颈、风险和状态。}
文章用 32 位可见深度排序加 16 位 tile 稳定排序减少 radix pass 和字节流量。
收益依赖 $P/M$、可见过滤成本和排序库。16 位 tile 键要求整个命名空间可表示，
多相机 tile 编码也必须算进去。本任务相机图像仅有 3,600 个 $16\times16$ tile，
但通用 API 不能因此永久假定 16 位足够。本 fork 仍走原有排序路径；本任务只做
源码评估，没有实现或测量该步骤，不能报告其增量收益。

\subsection{第五类：少做 SH、宽加载、base-2 指数与标量反向状态}

\subsubsection{可见性之后才读 SH}

三阶 SH 有 16 个基函数，每个 RGB 三通道，一共 48 个 FP32，即每高斯
192 字节。若先算完整 SH 再发现高斯在视锥外，就浪费了带宽和寄存器。
先做便宜的可见性判断，再读取 SH，是按数据依赖重排工作，不是数值近似。
连续系数用宽加载还能减少加载指令，但基函数与 RGB 的交错布局需要正确处理。

\begin{lstlisting}[language=C++,basicstyle=\ttfamily\small,breaklines=true]
Projection p = project_geometry(id);
if (!visible_and_has_tiles(p)) return;
float4 v0 = load_aligned_SH_group(id, 0);
float4 v1 = load_aligned_SH_group(id, 1);
evaluate_visible_color(v0, v1, /* more groups */);
\end{lstlisting}

\subsubsection{把自然指数改写为硬件 base-2 指数}

令 $\beta=-\tfrac12\log_2 e$，则
\begin{equation}
 \exp(-q/2)=2^{\beta q}.
\end{equation}
把 conic 的缩放提前到预处理，可以让内层直接形成 base-2 指数输入，使用硬件
快速指数路径。要注意交叉项 $2bd_xd_y$ 的系数约定：有的布局已把 2 合进去，
有的没有，拷贝三个数字而忽略该差异会渲染出错误椭圆。

\paragraph{瓶颈与风险。}
可见性门控减少 $M$ 对应的 SH 流量，宽读减少 LSU 发射压力，指数重排减少
内层转换工作。但是快速指数有近似和 FTZ 语义，系数预缩放也会改变舍入。
若原有 \texttt{\_\_expf} 加 fast-math 已经生成同样的硬件路径，仅把函数名改为
\texttt{exp2} 未必会少任何指令，不能从源代码名称推断收益。

\subsubsection{把反向的 RGB 状态投影成一个标量}

设像素的上游 RGB 梯度为 $\boldsymbol u=\partial L/\partial\boldsymbol I$，
当前高斯之后、以该处透射率归一的尾部颜色为 $\boldsymbol B_i$，其中包含背景。
在固定当前像素的 $\boldsymbol u$ 下，
\begin{equation}
 \frac{\partial L}{\partial\alpha_i}
 =T_{i-1}\left(\boldsymbol u\cdot\boldsymbol c_i
              -\boldsymbol u\cdot\boldsymbol B_i\right).
\end{equation}
后续计算只需要尾部颜色在 $\boldsymbol u$ 方向上的投影。因此令
$b_i=\boldsymbol u\cdot\boldsymbol B_i$，用一个标量递推取代 RGB 三个分量：
\begin{equation}
 b_{i-1}=\alpha_i(\boldsymbol u\cdot\boldsymbol c_i)
                  +(1-\alpha_i)b_i.
\end{equation}
这没有删掉颜色梯度本身：$\partial L/\partial\boldsymbol c_i
=T_{i-1}\alpha_i\boldsymbol u$ 仍然是向量。标量化的是反向遍历的运行状态。
上游梯度在像素内部固定、尾部递推线性，才允许这种投影交换。

\begin{lstlisting}[language=C++,basicstyle=\ttfamily\small,breaklines=true]
float behind_dot = dot(upstream_rgb, background_rgb);
for (splat : reverse_list) {
    float color_dot = dot(upstream_rgb, splat.rgb);
    float d_alpha = T_before * (color_dot - behind_dot);
    behind_dot = alpha * color_dot + (1 - alpha) * behind_dot;
    accumulate_color_and_geometry_gradients(d_alpha);
}
\end{lstlisting}

\paragraph{本分支与本任务。}
本 fork 的 SH VJP 已按颜色通道分工，不能把文章的 SH 寄存器模型原样套用。
本任务未新增交付 forward SH 宽加载或 base-2 重排；现有快速指数路径也不算新收益。
标量颜色投影及背景点积外提确实实现并测试过，112 个小例通过，但对应一轮
raster 中位数从 13.945 ms 变为 14.496 ms，没有收益，最终撤回。代数上减少
状态并不保证在特定编译器和 GPU 上更快；依赖链和寄存器分配同样重要。

\section{\texttt{s02}：让归约形状匹配梯度宽度}

\subsection{halving reduce-scatter 的核心不变量}

butterfly 计算了所有 lane 都不需要的重复结果。我们真正需要的是：每个分量在
\emph{某个} lane 上有完整总和，然后由那个 lane 写回。不需要 32 个 lane
都持有全部 9 项。这就是 reduce-scatter：在归约像素维度的同时，把输出分量分散。

用 4 个线程、4 个分量的教学模型说明。第一步让线程 $(0,2)$、$(1,3)$ 配对，
每对中的一个保留分量 0、1 的局部和，另一个保留 2、3。第二步在拥有同一分量区间
的线程间配对，继续把区间一分为二。最后每个线程只保留一个完整分量的和。

\begin{equation}
 \underbrace{4\times4}_{\text{每线程持有所有分量}}
 \ \longrightarrow\ \underbrace{4\times2}_{\text{合并两线程的部分和}}
 \ \longrightarrow\ \underbrace{4\times1}_{\text{每分量唯一拥有者}}.
\end{equation}

9 项不是 2 的幂，可以按实现的所有权规则补零、切分奇数区间，并保证虚拟项
不写回。文章给出的有效窗口长度为
\begin{equation}
 9\longrightarrow5\longrightarrow3\longrightarrow2
 \longrightarrow1\longrightarrow1,
 \qquad 5+3+2+1+1=12.
\end{equation}
因此从每 splat 45 条 shuffle 降至约 12 条，最后由 9 个拥有者在同一 warp
原子指令中写出 9 项。全局原子地址更新的总项数仍为 9，并没有神奇地合并成一个地址。

\begin{lstlisting}[language=C++,basicstyle=\ttfamily\small,breaklines=true]
// Conceptual collective, not a literal shuffle implementation.
Window w = initial_gradient_window(local_grad);
for (int stride = 16; stride != 0; stride >>= 1) {
    Ownership next = split_components_between_partners(w, lane, stride);
    Window remote = exchange_partner_selected_half(w, next, stride);
    w = add_selected_local_and_remote(w, remote, next);
}
if (owns_real_component(lane))
    atomicAdd(component_address(lane), owned_sum(w));
\end{lstlisting}

写真正的 shuffle 代码时，源 lane 要发送对方所需的分量，不是所有 lane 读取同一个
数组下标就会自动正确。最有效的验证是给每个输入 $(\ell,j)$ 唯一标记，在小模型中
追踪每一步的分量所有权，证明每个 $G_j$ 覆盖恰好 32 个输入且只写一次。

\subsection{两 splat half-warp folding：填满尾部的空闲位置}

9 个输出连半个 warp 都放得下。单 splat 的后期归约仍要跨 32 lane 合并，却已
不能再把一项拆成半项。文章将两个连续 splat A、B 的像素贡献同时保留：
先跨上下半 warp 配对，把 A 的 32 像素贡献折叠进低半 warp，把 B 折叠进高半 warp；
再在两个 16-lane 组内并行执行 reduce-scatter。

\begin{align}
 A'_{\ell j}&=A_{\ell j}+A_{\ell+16,j},\quad 0\leq\ell<16,\\
 B'_{\ell j}&=B_{\ell j}+B_{\ell-16,j},\quad16\leq\ell<32.
\end{align}
两组各自完成 16 像素局部和的再归约后，就得到各自完整的 32 像素结果。
初次 folding 交换 9 项，后面的半 warp 窗口归约花费 $5+3+2+1=11$，
合计每对 20 次，即平均每 splat 10 次，而不是单独执行两次 12。

\paragraph{预期瓶颈与风险。}
收益来自更少 shuffle 和更平行的写回，但代价是同时存活的两个 splat 状态。
寄存器数过高会降低驻留率，甚至 spill 到 local memory。文章的 SM80 启动约束
是其编译结果的一部分，不是跨 GPU 的最佳常数。还要处理奇数结尾、两 splat
不同的有效像素集合、absgrad 和非 RGB 通道。

\paragraph{本分支与本任务。}
最终只保留\textbf{单 splat halving reduce-scatter}，两 splat folding 未移植。
窄通道路径支持至多 4 个颜色/特征通道，并把可选绝对中心梯度纳入实际归约窗口；
通道更宽时沿用原实现。因此文章的 ``9 项、12 次''是教学例，不是本 fork 所有配置
的指令计数。

一次后续检查发现 $4\times4$ tile 只有 16 个线程：假定完整 warp 的拥有者可能
根本不存在，导致本应非零的不透明度梯度变零。\textbf{传入 active mask 不能创造
缺失的拥有者。}最终提交 \texttt{88513ea} 对
\texttt{block\_size \% 32 != 0} 的整个 block 回退到旧路径，覆盖 $4\times4$
及 $6\times6$ 等配置。最终 RGB backward 报告 48 个寄存器、无 stack/local
memory；早期 scatter-only 的静态指令计数不能冒充含回退分支的最终二进制总计数。

\section{\texttt{s03}：SH 寄存器流式化与协作清零}

\subsection{为什么 SH 不需要同时保存全部输入和输出}

设方向为 $\boldsymbol d$，SH 基函数为 $Y_k$，RGB 系数为 $\boldsymbol s_k$：
\begin{align}
 \boldsymbol c(\boldsymbol d)&=\sum_kY_k(\boldsymbol d)\boldsymbol s_k,\\
 \frac{\partial L}{\partial\boldsymbol s_k}
 &=Y_k(\boldsymbol d)\boldsymbol u,\\
 \frac{\partial L}{\partial\boldsymbol d}
 &=\sum_k(\boldsymbol u\cdot\boldsymbol s_k)
                 \nabla_{\boldsymbol d}Y_k(\boldsymbol d).
\end{align}
系数梯度只依赖该基函数和上游梯度，写完就可以释放寄存器；方向梯度只需要一个
三分量运行和。于是没有必要同时保留 48 个输入系数和 48 个输出梯度。
颜色非负截断等前向操作的导数掩码应先作用于 $\boldsymbol u$，不能在流式化时丢掉。

文章把 16 个三阶基函数分成 4 组，每组 4 个，即 12 个 RGB 浮点数。
12 同时能被 RGB 宽度 3 和向量读取宽度 4 整除，适合作为一个寄存器窗口。

\begin{lstlisting}[language=C++,basicstyle=\ttfamily\small,breaklines=true]
float3 d_direction = make_float3(0, 0, 0);
for (int group = 0; group < 4; ++group) {
    float v[12];
    load_three_float4(group, v);
    for (int k = 0; k < 4; ++k) {
        Basis y = basis_and_direction_derivative(4 * group + k);
        store_SH_gradient(4 * group + k, y.value * upstream_rgb);
        d_direction += dot(upstream_rgb, rgb(v, k)) * y.derivative;
    }
}
propagate_normalized_direction_vjp(d_direction);
\end{lstlisting}

\paragraph{预期瓶颈与风险。}
这针对活跃寄存器数量和由此产生的驻留率，而非减少 SH 的数学项数。
文章的 FP32 backward 寄存器数从 152 降至 104，是 A100 具体内核的结果。
窗口太小会增加加载与索引，太大又失去驻留率收益。方向归一化的雅可比仍需正确
回传，不能把方向导数直接当作位置导数。

\subsection{移除 bulk memset 不是允许输出未初始化}

若完整参数梯度缓冲区先清零，再由可见高斯覆盖，每个活跃元素被写了两次。
文章让活跃项直接写最终梯度，非活跃项由内核写零，移除单独的全量 memset。
但有一个所有权前提：这些输出必须是最终赋值，而不是依赖预先清零的原子累加目标。
光栅阶段共享的累加梯度缓冲区不能不加区分地照搬该方法。

朴素地让每个线程清自己高斯的 192 字节 SH 行，会使 warp 同一条存储指令的地址
相隔 192 字节。A100 的缓存按 32 字节 sector 传输，许多 sector 只写一小部分。
更好的方式是让 warp 协作线性扫描 32 个高斯的连续 SH 区域，以 16 字节块为单位：

\begin{lstlisting}[language=C++,basicstyle=\ttfamily\small,breaklines=true]
unsigned active = __ballot_sync(0xffffffffu, gaussian_is_active);
// 32 rows, each with 12 float4 chunks for degree-3 RGB SH.
for (int chunk = lane; chunk < 32 * 12; chunk += 32) {
    int owner = chunk / 12;
    if (((active >> owner) & 1) == 0)
        store_float4_zero(warp_SH_base + chunk);
}
// Active rows are fully assigned by their VJP owners.
\end{lstlisting}

相邻 lane 总访问相邻 16 字节块，因此合并性改善，而非活跃行仍有正确的零梯度。
尾 warp、实际系数数 $K$ 大于当前阶数所用数量、未启用方向梯度等情况，都需要
定义哪些元素由 VJP 写、哪些由零填充写，确保没有遗漏和重叠竞态。

\paragraph{本分支与本任务。}
本 fork 的 SH backward 是 $N\times3$ 个线程，每线程处理一个颜色通道，
在 VJP 中直接写系数梯度，不存在文章同样的 ``48 输入加 48 输出都常驻''结构。
数据又是一阶 SH，不能套用三阶的带宽与寄存器收益。
本任务实现过选择性协作清零，288 个 SH 回归例覆盖阶数 0--4、mask、超额 $K$
和方向梯度开关，全部通过；但配对完整渲染测量没有可靠收益，一轮中位数为
64.056 ms 对 64.973 ms，最终撤回。\texttt{zeros\_like} 改成
\texttt{empty\_like} 而不负责所有未写元素，绝不是这一优化的正确简化。

\section{\texttt{s04}：前向知道谁贡献过，就不要让反向重新猜}

\subsection{几何重叠与真实贡献的区别}

一个高斯可以与 tile 椭圆相交，却因为低权重或遮挡没有对任何像素产生实际混合。
几何 mask 是保守候选集合；贡献 mask 是执行了基线接受逻辑之后得到的事实。
记 tile 中第 $i$ 个高斯的 warp 贡献掩码为
\begin{equation}
 m_i=\sum_{w=0}^{7}2^w
 \mathbf{1}\{\text{warp }w\text{ 至少一个像素实际混合了 }i\}.
\end{equation}
当 $m_i=0$，该高斯在这个 tile 内可以从反向列表中删去。注意这里要求的是
``实际混合''，不是仅满足 $\alpha\geq\tau$；已经终止的像素不能使它成为贡献者。

\subsection{原位压缩和停止索引的重映射}

设原列表中的贡献标记为 $v_i=\mathbf{1}\{m_i\ne0\}$，稳定压缩后的位置为
\begin{equation}
 r_i=\sum_{j<i}v_j.
\end{equation}
保留原顺序，把第 $i$ 项移到 $r_i$，于是压缩后的顺序仍是深度顺序。
一个像素原来的停止位置也必须用相同前缀计数映射；否则反向可能在错误的 splat
停止，甚至使用一个压缩前位置索引压缩后数组。

\begin{lstlisting}[language=C++,basicstyle=\ttfamily\small,breaklines=true]
stage_whole_batch_before_overwriting_source();
for (splat : batch) {
    unsigned cmask = collect_actual_blend_bits_from_all_warps(splat);
    keep[splat] = (cmask != 0);
}
rank = exclusive_scan(keep);
for (splat : batch) {
    if (keep[splat])
        list[tile_begin + prior_kept + rank[splat]] = pack(splat);
}
remap_per_pixel_last_indices_using_same_prefix();
\end{lstlisting}

``原位''的安全依据是输出写指针不超过当前读指针，并且当前批次已完整 staging。
若某 warp 提前覆盖另一个 warp 尚未读取的元素，数学上的稳定压缩依然会在实现中
变成数据竞争。tile 的切片、跨 batch 的计数和同步点都必须纳入证明。

\subsection{高字节 \texttt{cmask} 与全局 \texttt{touched}}

文章将 8 位贡献掩码打包进 32 位列表项的高字节：
\begin{equation}
 \text{packed}=\left(m_i\ll24\right)\mathbin{|}\text{id},\qquad
 \text{id}=\text{packed}\mathbin{\&}(2^{24}-1).
\end{equation}
反向直接使用该掩码，避免重新计算子块几何边界。另设每个高斯的
\texttt{splat\_touched} 标记，只有实际贡献过的高斯才运行昂贵的预处理 VJP。
这会把``视锥可见但完全遮挡''的高斯从 SH backward 工作集中移除。

\paragraph{24 位限制不能靠场景平均规模回避。}
这种格式只容纳 $0$ 至 $2^{24}-1$ 的 ID。本任务有 20,789,814 个高斯，
超过 $2^{24}=16,777,216$；最终实测列表 ID 达到 20,789,812。
直接截低 24 位会把某些高斯别名到另一些高斯，输出可能仍然看似合理，却产生
错误梯度。安全设计可以采用独立 mask 缓冲区、64 位条目、可见集合重映射，
或明确的按规模回退；每一种都有额外字节、扫描和接口代价。

\paragraph{预期瓶颈与风险。}
前向额外做投票、计数和写回，以换取反向少 staging、少求值、少 VJP。
因此纯前向可能更慢。\texttt{touched} 还需要正确初始化、跨 tile 更新与跨内核
可见性；不能仅因为多个线程都写 1 就忽略内存模型或竞争检查。
多通道分块调用、多个消费者复用原列表时，修改输入列表尤其危险。

\subsection{本任务选择更小的改动：只裁掉完全无用的尾部}

完整原位压缩、打包 \texttt{cmask} 和新的 \texttt{touched} 契约都没有移植。
本任务利用已有逐像素 \texttt{last\_ids} 求 tile 内最大停止索引，只删去所有
像素都未使用的后缀：
\begin{equation}
 L_{\mathrm{tile}}=\max_{p\in\mathrm{tile}}L_p.
\end{equation}
反向从这个边界开始，不再 staging 更深的完全遮挡后缀；中间没有贡献的洞仍保留。
该方案不修改共享列表，不压缩 ID，也不改变停止索引的坐标系，适合本 fork 的
通道分块复用。它不是文章完整压缩的同义词。

空 tile 必须先统一返回；从包含端点变成排他端点的加一应先提升整数宽度，避免
靠近整数上界时溢出。初期隔离实验中，在 scatter 之上加入后缀裁剪有明显收益，
因而最终保留；但最终总倍率来自重建后二进制的重新配对测量，不能把早期倍率
与后期倍率相乘来宣称最终性能。

\section{\texttt{s05}：分歧、谓词化与精确的数值分支边界}

\subsection{为什么 ``少执行线程''未必意味着少发射指令}

warp 中各 lane 共享指令发射。若 31 个像素都 miss、1 个像素 hit，普通分支下
warp 仍要发射混合和梯度指令，只是其中多数 lane 被禁用。若候选 mask 已经保证
至少一个像素有用，逐像素的多个 \texttt{continue} 常常只增加重汇合与谓词管理。
让 miss lane 的有效 $\alpha=0$，可以把递推写成连续直线代码：
\begin{equation}
 \alpha_{\mathrm{eff}}=
 \begin{cases}\alpha,&\text{hit},\\0,&\text{miss},\end{cases}
 \qquad
 \Delta\boldsymbol C=T\alpha_{\mathrm{eff}}\boldsymbol c,
 \qquad T'=T(1-\alpha_{\mathrm{eff}}).
\end{equation}
在有限数、没有副作用的理想代数中，miss lane 的状态完全不变。

\begin{lstlisting}[language=C++,basicstyle=\ttfamily\small,breaklines=true]
float power = exact_shared_quadratic(conic, dx, dy);
float raw = opacity * fast_exp(power);
float capped = min(reference_cap, raw);
bool hit = active && reference_power_test(power)
                  && reference_cutoff_test(capped);
float alpha = hit ? capped : 0.0f;
// Reference termination ordering must still be applied here.
C += (T * alpha) * color;
T *= (1.0f - alpha);
\end{lstlisting}

\paragraph{本 fork 没有文章同样的先决条件。}
本任务未交付精确 warp 贡献 mask，所以候选高斯仍可能整 warp 都 miss。
保留 warp-wide all-miss 跳过才合理；不能借文章的结构不变量，删除本实现中
仍有价值的统一分支。宽通道路径也保留原来的条件求值，避免所有 lane 对大量
特征执行无效运算。

\subsection{精确舍入：前后向必须在同一侧跨过 cutoff}

若前向把某像素判为 hit，而反向因 FMA 融合方式不同判为 miss，反向会漏掉一次
$T/(1-\alpha)$ 恢复，后面整条反向递推都会受影响。问题不只是一个高斯边缘上
极小的局部误差，而是离散控制流变化。

因此文章用显式 round-to-nearest-even 操作固定二次型求值的表达式树。
下面只是展示固定树的方法；实际树必须与目标基线一致：

\begin{lstlisting}[language=C++,basicstyle=\ttfamily\small,breaklines=true]
// Illustrative RN tree; use exactly the same helper in fwd and bwd.
float xx = __fmul_rn(dx, dx);
float xy = __fmul_rn(dx, dy);
float yy = __fmul_rn(dy, dy);
float axx = __fmul_rn(a, xx);
float bxy = __fmul_rn(two_b, xy);
float cyy = __fmul_rn(c, yy);
float q = __fadd_rn(__fadd_rn(axx, bxy), cyy);
float power = __fmul_rn(-0.5f, q);
\end{lstlisting}

不能把这个教学树直接覆盖到已有 fork：若原基线使用另一种乘加树，固定为新树
虽然让新前后向一致，仍可能改变与旧实现的 cutoff 归属。本任务没有交付新的
前向表达式、conic 预缩放或 base-2 重排，而是保持现有前向语义进行差分验证。
所谓 ``exact rounding'' 是操作舍入可控，不是高斯函数的实数精确求值。

\subsection{clamp、cutoff 和反向零贡献不是同一件事}

cutoff 是``不接受这个混合项''；cap 是``接受，但把 alpha 饱和''。
对于 cap 已饱和的像素，颜色梯度仍然存在，而通过 alpha 到几何、不透明度的导数
应按参考实现的饱和导数规则门控。不能把这两类掩码合成一个随意的 \texttt{hit}。

另外 IEEE 浮点满足 $0\times\infty=\mathrm{NaN}$，并不满足
``乘了零就是零''。有限 conic 也可能在中间乘法溢出；被拒绝 lane 的上游梯度
还可能是非有限值。若先无条件算出 NaN，再乘零，归约就会污染其他像素。

\begin{lstlisting}[language=C++,basicstyle=\ttfamily\small,breaklines=true]
Grad g = exact_zero_grad();
if (hit) {
    update_backward_color_state();
    g.color = accepted_color_vjp();
    if (reference_alpha_derivative_is_enabled(raw)) {
        g.geometry = geometry_vjp();
        g.opacity = opacity_vjp();
        g.abs_mean = optional_abs_mean_vjp();
    }
}
reduce_gradients(g);  // invalid lanes really contribute zero
\end{lstlisting}

这说明安全实现可以只让某些递推谓词化，仍对危险导数求值保留显式门控。
追求源代码 ``没有 if'' 不是目标；减少真实机器码分歧并保留语义才是目标。
NaN 的比较、最小值指令行为、clamp 边界的等号约定也要继承基线。

\paragraph{本任务结果。}
窄通道 backward predication 在加固后保留，前向完全不变。加固内容包括拒绝 lane
不更新颜色运行缓冲，拒绝或饱和 lane 不计算几何、不透明度及 absgrad 导数。
一次最终集成消融中，相机 0 在 scatter 加 trim 之上约有 $1.038\times$ 的
增量 raster 收益；这是该测量序列的结果，不是任意场景或最终全部相机的独立倍率。
有限但溢出的 conic、非有限拒绝 lane 和非默认流已进入最终回归矩阵。

\section{\texttt{s06}：多行映射与执行流水线平衡}

\subsection{每线程两个竖直像素：先在线程内合并}

文章将一个 warp 的像素覆盖从 $8\times4$ 改为 $8\times8$，每个线程负责
同一列的两个相邻像素。一个 $16\times16$ tile 从 256 线程变为 128 线程，
即由 8 个 warp 变为 4 个 warp。两个像素共享 $d_x$、横向二次项和高斯属性设置。
反向先在线程寄存器中把两个像素的梯度相加，再做 warp 归约，故每 tile 的归约
次数和对应原子小计数减半。

\begin{lstlisting}[language=C++,basicstyle=\ttfamily\small,breaklines=true]
Pixel top = pair_top_pixel(lane, warp_id);
Pixel bottom = top + vertical_offset(1);
Column col = setup_shared_x_terms(splat, top.x);
State s0 = update_one_row(col, top, state0);
State s1 = update_one_row(col, bottom, state1);
Grad local = s0.grad + s1.grad;
warp_reduce_scatter(local);
\end{lstlisting}

两个像素仍有独立颜色、透射率、上游梯度和终止状态，不能因为共享列就共用一份
像素状态。屏幕底边和不整齐图像边缘可能只有其中一个像素有效。

\subsection{warp 私有 staging：去掉等待最慢邻居的时间}

旧结构的所有 warp 共用一个 tile 级共享 batch；即使某 warp 的区域很空，也必须
在 \texttt{\_\_syncthreads()} 等待最繁忙的 warp。文章在 backward 给每个 warp
分配独立的 32-splat staging 切片，只在本 warp 内同步，并根据贡献 mask 独立遍历。
这不是把同步函数名直接改成 \texttt{\_\_syncwarp()}：共享内存的所有权和生产消费
范围必须一起重设计，否则其他 warp 可能覆盖仍在读取的数据。

\paragraph{流水线视角。}
文章用合成场景的 Nsight Compute 指标说明，某些步骤受 block barrier 拖累，
另一些步骤受 FMA、ALU 或总发射率限制。各执行管线都不到满载，并不表示没有瓶颈：
如果 scheduler 每周期只能发射有限条指令，混合整数、浮点、指数和加载指令的总量
可以先耗尽发射带宽。要区分 ``少 FMA''和``少总指令''，也要区分运行时动态计数
与静态 SASS 中出现了多少条指令。

\subsection{把不透明度并入指数，复用列常数}

对 $o>0$，有
\begin{equation}
 o\,2^p=2^{p+\log_2o}.
\end{equation}
预处理保存 $\log_2o$，列设置时将它合进横向常数，每行指数输出就直接是
$\alpha_{\mathrm{raw}}$，少一次逐行不透明度乘法。反向几何使用的 log-alpha
导数也可保留 $o$ 因子，到输出不透明度梯度的拥有者处再恢复相应比例。

代数相同不等于浮点位相同：$o=0$、极小 $o$、FTZ、先乘后 cap 与先加 log
再指数，都会涉及边界。文章还结合列范围检查减少逐像素 power 判断，这依赖其
完整设计，不是只替换一个指数就能继承的性质。

\subsection{竖直两行矩：把乘法延后到归约之后}

取文章的符号约定，两个像素的相对纵坐标为 $y$ 与 $y-1$，相应 log-alpha
梯度为 $g_0,g_1$。定义三个矩：
\begin{align}
 m_0&=g_0+g_1,\\
 S_y&=g_0y+g_1(y-1)=m_0y-g_1,\\
 S_{yy}&=g_0y^2+g_1(y-1)^2=m_0y^2-g_1(2y-1).
\end{align}
这些量足以构造多种均值和 conic 梯度。如果竖直 shuffle 连接的是同一像素列，
则各线程有同一个 $d_x$，可以利用
\begin{equation}
 \sum_\ell d_x m_{0,\ell}=d_x\sum_\ell m_{0,\ell}
\end{equation}
把横向乘法延后到竖直部分归约之后。这里节约来自几何映射提供的公共因子。
若修改 lane 编号使配对不再同列，上式的公共 $d_x$ 假设就失效了。

\subsection{终止状态编码与 BF16 的不同性质}

文章还用透射率的符号位编码像素是否已终止，避免单独维护某些布尔状态：
负号表示完成，绝对值仍保存输出所需透射率。这只适用于 $T$ 原本非负的约定；
所有使用点都得区分``存储状态''和``物理透射率''，背景混合或恢复导数时不能
误用负值。负零、下溢和最终取绝对值也需要明确语义。

BF16 SH 则属于另一类优化：通过降低系数存储宽度减少带宽，而不是单纯重排计算。
BF16 保留与 FP32 相同数量的指数位，但有效精度更低；应说明是否以 FP32
累加、梯度如何保存、优化器是否有主副本，以及参数更新时何处发生量化。
文章在其测试中报告相同到 0.01 dB 的 PSNR，不等于任意场景的逐元素结果相同，
更不能自动证明任意优化器的训练轨迹不变。

\paragraph{预期瓶颈、风险和本任务状态。}
多行减少设置、归约和原子次数，warp 私有 staging 减少 barrier 等待；代价是
每线程更多像素状态和更复杂 mask 映射。RGB+ED 或宽特征增加寄存器压力，可能
使 RGB 上的收益反转。BF16 对三阶 SH 更有动力，本任务一阶系数流量本来就小。
上述多行、私有 staging、管线数学重排、符号编码和 BF16 都只做了源码/硬件
适用性评估，\textbf{未实现、未测量、未作为交付内容}。不能给这些步骤填入
来自 A100 的预期倍率，冒充 4090 上已经获得的收益。

\section{负结果：失败实验教会了什么}

文章单列了 17 个负结果及硬件陷阱。它们的教学价值不是形成一份永远禁止使用的
技巧名单，而是说明同一个局部优化可以撞到另一个约束。下面按因果关系整理这些
陷阱，并与本任务实际撤回的实验区分；不把未在本任务运行的方案说成已经失败。

\subsection{驻留率不是越高越好，寄存器也不是越少越好}

强制压低寄存器可以增加理论驻留 warp 数，却可能导致 spill，产生 local memory
读写；反过来，放宽寄存器预算可以消除 spill，却可能跨过驻留 block 的台阶。
正确的问题不是 ``occupancy 达到多少''，而是新增的 warp 是否能隐藏当前瓶颈的
延迟，以及为它付出了多少额外流量。一个依赖链很长的指数--除法递推，也不会仅靠
少两条独立加法就完全变快。

因此应该联合记录寄存器、spill、活跃 warp、内核时间，不能只展示源代码行数减少。
本任务最终保留的 reduce-scatter 需要更多局部部分和，但仍无 local memory；
两 splat folding 则没有测量，不能假设进一步减少 shuffle 必然更快。

\subsection{共享内存缓存并不免费，更多并行也可能引入等待}

把重复读取放进共享内存可能节约全局加载，也增加 staging、屏障、bank conflict
和容量占用。将工作拆成更多 warp 可能改善覆盖，却增加归约次数；把工作合并给
更少线程可能降低设置成本，却增加寄存器和长尾。文章的方形多行映射与私有 staging
是配套变化，不是两个可以无条件相加的独立倍率。

\begin{lstlisting}[language=C++,basicstyle=\ttfamily\small,breaklines=true]
// A useful experiment record contains both intent and a rejection gate.
Hypothesis h = "less staging or fewer issued instructions";
Candidate c = change_one_mechanism(h);
if (!differential_correctness(c)) reject(c);
else if (!paired_repeated_target_timing_improves(c)) reject(c);
else retain_with_documented_scope(c);
\end{lstlisting}

\subsection{缓存局部性和排序的收益会互相影响}

宽键变窄会少搬字节，但新增过滤、扫描和中间数组也有成本；tile 更精确会减少
排序规模，但也可能让其他固定开销占比升高。重新排列高斯以改善某阶段的连续读，
可能让另一个阶段的梯度写回更分散。L2 命中率单独变高并不是充分证据：如果总请求
数和指令数一起增长，绝对时间仍可能更差。

文章的两阶段排序针对重复深度排序这一具体问题。本任务没有实现它，所以本任务
的 ``未交付''不能被写成 ``实测无效''。同样，AABB 子块原型的弱收益也不能用于
否定文章精确 conic mask 加紧凑布局的组合。

\subsection{本任务三个明确撤回的实验}

\begin{enumerate}
 \item \textbf{标量颜色累加。}数学线性化成立，112 个小回归通过，但对应
 raster 测量没有收益，故恢复原计算。这里的结论是目标负载上不值得交付，
 不是标量投影公式错误。
 \item \textbf{SH 选择性协作清零。}288 个专门测试通过，但完整渲染配对结果
 没有可靠增益。索引、除法、掩码开销和非活跃行比例可能抵消减少的写流量；
 这是解释假说，而不是已经由硬件计数器证明的瓶颈诊断。
 \item \textbf{前向 AABB warp culling。}通过小回归但收益小、噪声显著，
 专用循环和病态边界回退增加维护负担。早期 FP64 边界原型还暴露了与浮点溢出
 基线兼容的风险，所以没有交付任何前向源码修改。
\end{enumerate}

此外，部分 warp 缺失拥有者不是一次可以接受的性能负结果，而是必须修复的
正确性缺陷。最终引入回退并重新构建、重新测量，先前二进制退出交付，不能继续
使用它更漂亮的历史数字。判定性能之前必须先确认测量的就是待交付二进制。

\section{把文章读完之后：怎样评价这次移植}

\subsection{交付范围与最终测量}

最终提交 \texttt{88513ea} 的内核变化仅在 backward rasterization 和归约辅助
代码：窄通道 reduce-scatter、遮挡后缀裁剪、安全的 alpha predication，以及
部分 warp 回退。前向和 SH 源码已恢复，完整压缩、两阶段排序、two-splat folding、
多行和 BF16 均未移植。该状态是隔离的优化/审查分支，不是已经替换用户安装包
或合并到主工作区的生产发布。

最终重建后的全量场景 RGB 配对结果如下。raster 指前后向光栅 API；完整渲染
还包括投影、SH、分箱和相应反向等。数字是作业中位数的中位数，不是取最快一次。

\begin{center}
\begin{tabular}{rcc}
\hline
相机索引 & raster 倍率 & 完整可微渲染倍率 \\
\hline
0    & $1.337\times$ & $1.100\times$ \\
2552 & $1.648\times$ & $1.016\times$ \\
5103 & $1.363\times$ & $1.039\times$ \\
\hline
\end{tabular}
\end{center}

这只有三个视图，不是全部 5,104 个相机。相机 2552 的两个完整渲染配对方向
不一致，其小幅汇总收益仍不确定。产品 RGB+ED 加 absgrad 设置在相机 0 上为
$1.311\times$ raster、$1.081\times$ 完整渲染；不能用 RGB 三通道结果代替
四通道产品设置。未测试优化器、损失、训练循环整体收益或训练收敛。

\subsection{正确性证据也需要分层}

最终二进制在默认流和非默认流上各通过 136 个差分测试，memcheck 和 synccheck
均报告 0 错误。全量场景 RGB/alpha 与基线逐位相同，直接屏幕空间梯度的最差
relative-L2 为 $2.62\times10^{-6}$（RGB）及 $1.48\times10^{-6}$
（RGB+ED），低于保持不变的 $10^{-4}$ 门槛。

但是完整参数路径的单次 raw-quaternion 梯度比较约为 $1.63\times10^{-4}$，
\textbf{没有通过严格的 $10^{-4}$ 门槛}。重复原子累加噪声和近退化投影会放大
误差，是理解现象的上下文，而不是豁免阈值的理由。固定 16 次运行、CPU FP64
累积的平均梯度检查通过，四元数平均 relative-L2 分别为
$4.57\times10^{-5}$ 和 $3.50\times10^{-5}$；这只支持均值一致性，
不把单次失败变成通过，更不构成收敛证明。

\subsection{一套可以带走的思考顺序}

\begin{enumerate}
 \item \textbf{固定语义。}先写清参数激活、像素中心、cutoff、cap、终止顺序、
 相等深度、空输入与流语义。否则所谓加速可能只是换了问题。
 \item \textbf{减少集合。}先问能否减少 $P$、实际遍历次数和不必要 SH 行，
 再问单次求值能否少一条指令。
 \item \textbf{匹配硬件形状。}把 9 项输出映射到拥有者，把 SH 活跃窗口限制在
 必要范围，把共享数据与线程像素映射一起设计，而不是照抄某个函数名。
 \item \textbf{检查代数背后的浮点。}线性投影、log 指数合并和零 alpha 都有
 有效范围；边界比较、非有限中间值及原子顺序仍然决定实现正确性。
 \item \textbf{逐步过门。}每个改动独立建立假说、差分回归、消融和配对时间；
 拒绝无收益方案，保留不确定性，不把未实现部分写成已完成结果。
\end{enumerate}

文章最值得迁移的不是 $7.30\times$ 这个数字，而是这条从算法工作量到硬件映射、
从数值契约到证据边界的推理链。对本 fork 而言，较小且可验证的 backward 改动
已经带来了可重复的 raster 收益；剩余方案应作为新的、独立的研究问题，
而不是从 A100 图表中预支的性能承诺。
